% ===== 原文 pp.2（下半）–4（上半） =====
% 本块由 Lead 负责：§1 Introduction 全文（含脚注 1、2）+ 式 (1.1)；tag 区间 1.1

% [原文 p.2]
\section{Introduction}

本文的目的，是对 2d/4d 对应关系给出一个初等的介绍；该对应关系最初发现于文献 [2,3]。
让我们先给出其核心思想——它其实相当简单。

% [Tutor Note: p.2–3, §1 开篇。这段「六维理论在两个流形上两种顺序求配分函数」的论述是全篇
%  的骨架，建议先把它记成一句话：$\cS(X_4\times C_2)$ 有两种「先缩哪一个」的算法，一种给出
%  四维理论在 $X_4$ 上的配分函数，另一种给出二维理论在 $C_2$ 上的配分函数；2d/4d 对应关系
%  就是这两种算法给出同一个数。后文 §5.1、§5.2 不过是把这个「同一个数」在两个具体背景
%  （$X_4=S^1\times S^3$ 与 $X_4=S^4$）上算出来而已。]

我们从一个六维量子场论 $\cS$ 出发，考虑它在乘积流形 $X_4\times C_2$ 上的配分函数，其中
$X_4$ 是四维的、$C_2$ 是二维的。进一步假设：由于超对称扭转，所得配分函数只依赖于 $X_4$ 与
$C_2$ 的\textbf{形状}而不依赖于它们的\textbf{大小}。于是六维配分函数 $\cS(X_4\times C_2)$ 可以用两种方式
计算。一方面，若把 $C_2$ 取得很小，我们先得到一个四维理论 $\cS(C_2)$，再考虑它在 $X_4$ 上的
配分函数，即 $\cS(C_2)(X_4)$。另一方面，若把 $X_4$ 取得很小，我们先得到一个二维理论
$\cS(X_4)$，再考虑它在 $C_2$ 上的配分函数，即 $\cS(X_4)(C_2)$。于是我们得到一个等式

\begin{equation}
  \cS(X_4)(C_2) = \cS(C_2)(X_4). \tag{1.1}
\end{equation}

% [Tutor Note: p.3, 式 (1.1)。注意「形状不变、大小不变」这一条假设（shape-dependence 而非
%  size-dependence）是等式成立的全部前提：它来自超对称扭转（supersymmetric twist），
%  §4.1 会解释为什么配分函数可以只依赖形状。若没有这一条，(1.1) 左右两边会分别依赖
%  $X_4$ 与 $C_2$ 的体积，等式就不成立。]

到此为止，构造是极其一般的。要得到具体的东西，我们必须做选择。举例来说，取 $\cS$ 为
$A_1$ 型的六维 $\cN=(2,0)$ 理论，$C_2$ 为任意一个二维曲面。由此得到的四维理论
$\cS(C_2)$ 通常称为 $A_1$ 型的 class S 理论，它是一类 $\cN=2$ 超对称理论：若干 $\SU(2)$
规范群因子，耦合到若干三基表示场（trifundamental field）[4]。

% [Tutor Note: p.3。$A_1$ 型的 6d $\cN=(2,0)$ 理论其规范代数就是 $\SU(2)$；之所以写成
%  $A_1$ 是沿用 ADE 分类的记号（$A_1=\fsu(2)$）。后文 §5.3.1、§6.3 会把这一记号推广到
%  $A_{N-1}$ 等其它类型，那时 $T_N$ 理论登场。]

若取 $X_4=S^4$，配分函数 $\cS(C_2)(X_4)$ 可以用局域化来计算 [5,6]；从计算结果可以看出，
$\cS(S^4)$ 就是 Liouville 理论，它是一个非紧致的二维共形场理论 [2]。若取
$X_4=S^3\times S^1$，配分函数 $\cS(C_2)(X_4)$ 称为超共形指标（superconformal index）
[7,8]；从计算结果可以看出，$\cS(S^3\times S^1)$ 就是 $q$-形变的 $\SU(2)$ 杨--米尔斯理论
[3,9]。这些是对应关系最简单的情形；各种推广都是可能的，而且已经被做出来了。

% [Tutor Note: p.3。这两句是全篇的「结论预告」，务请对照着读：$S^4$ 局域化 $\leftrightarrow$
%  Liouville 理论（§5.2 给出的等式是 $S(S^4)=\Liouville$ 理论），$S^1\times S^3$ 上的
%  （超共形）指标 $\leftrightarrow$ $q$-形变 $\SU(2)$ 杨--米尔斯（§5.1 给出的等式是
%  $S(S^1\times S^3)=$ 二维 $q$-形变 $\SU(2)$ 杨--米尔斯）。注意本文里「$S^3\times S^1$」
%  与「$S^1\times S^3$」是同一个空间，作者两种顺序混用（沿用引用文献的写法），
%  不代表不同对象。]

% [Tutor Note: p.3。原文此处说 "$\cS(S^4)$ is the Liouville theory"、"
%  $\cS(S^3\times S^1)$ is the $q$-deformed SU(2) Yang-Mills theory"：严格讲是指
%  「依赖 $\cS$ 的那个二维理论的配分函数」与「相应二维理论的配分函数」相等（即下面的
%  (5.25)、(5.26)）；这里先给出简写，读 §5 时不要被简写误导。]

这些对应关系最初被发现至今已经五年了\footnote{本文的主体写于 2014 年 8 月，仅于 2016 年
夏对参考文献作了小幅更新。}，而数不清篇幅的原始论文已经被耗费在描述其中的细节与各种推广上。
专注于此主题的综述文章数目也已经相当庞大。然而，如此巨量的信息也会令人望而生畏；
作者因此觉得，再写一篇简明的综述并非完全没有用处——这样一位新产生兴趣的读者就能快速通读，
对「这一对应关系是如何被发现的、细节可以到哪里去学、还有哪些尚未解决的问题值得研究」
形成一个大致概念。

% [Tutor Note: p.3。原文用词是 "countless pages of original articles have been already
%  wasted to describe the details"——"wasted"（浪费）是本底本的原文措辞（文字层与页图一致，
%  非抽取错误）。译文照原文译出，不作径改；读者若觉得语气偏重，可理解为作者对文献膨胀的
%  自嘲，不宜据此推定此处必有笔误。]

因而本文有意写成一份浅显的概览。许多事实将作为事实直接陈述，而不加解释。细节与微妙之处
会被提及，但不会完整处理；取而代之，我们会给出指向综述或原始论文的文献。\footnote{这些
参考文献并不详尽，甚至连广泛都谈不上。作者很乐意在 arXiv 版本中加入更多条目，故请不吝来信告知。}
叙述也不会是完全合乎逻辑的。最系统的做法是从六维出发，仔细分析紧致化，最后一步才抵达
这一对应关系。得益于近来的进展，按这个顺序写一篇综述并非不可能；但那样对初入此领域的人
会是一份极难读的材料。

% [Tutor Note: p.3, 脚注 2 的位置。作者坦承「引用并不详尽、甚至连广泛都谈不上」，并请读者
%  直接给他发邮件补充。阅读本文时，凡看到「details can be found in ...」而正文未展开的地方，
%  都应当查它所引的那一篇；本文的定位是地图，不是教科书。]

因此，本文的组织方式是解释这一对应关系\textbf{如何}运作，而不是\textbf{为什么}会有这一对应关系。
在 §2 中，我们先学习两个对我们很重要的二维场理论：二维 $q$-形变杨--米尔斯理论与
Liouville 理论。在 §3 中，我们直接以四维场理论的方式引入 $\SU(2)$ 型的 class S 理论——
它们由与 Riemann 曲面（以及该曲面的三穿刺球面分解）相关联的 Lagrangian 定义。在 §4 中，
我们快速引入超对称局域化技术，并说明 $S^1\times S^3$ 上或 $S^4$ 上的配分函数如何计算。
在 §5 中，我们把超对称局域化用于 $\SU(2)$ 型的 class S 理论，将会看到 $S^3\times S^1$ 上
与 $S^4$ 上的配分函数分别由 $q$-形变杨--米尔斯理论与 Liouville 理论给出。随后我们解释
这一对应关系如何能从六维 $\cN=(2,0)$ 理论的角度来理解。最后，我们在 §6 中回顾可能的
未来方向。

% [Tutor Note: p.3–4, §1 末段。这段是全书阅读指南，等价于一份「章节依赖图」：
%  §2（两个二维理论）+ §3（class S 四维理论）+ §4（两种超对称背景与局域化）
%  → §5（两个对应关系 + 六维解释）→ §6（开放问题）。
%  若只想快速抓住对应关系本身，最短路径是：§2.1.3–2.1.4（$q$-形变杨--米尔斯）+ §2.2.3（Liouville
%  的共形块结构）+ §4.2–4.3（两个背景上的局域化）+ §5.1–5.2（两个等式），其余可后补。]
